% **************************************************************************************************************
% A Classic Thesis Style
% An Homage to The Elements of Typographic Style
%
% Copyright (C) 2015 André Miede http://www.miede.de
%
% If you like the style then I would appreciate a postcard. My address 
% can be found in the file ClassicThesis.pdf. A collection of the 
% postcards I received so far is available online at 
% http://postcards.miede.de
%
% **************************************************************************************************************
\RequirePackage{fix-cm} % fix some latex issues see: http://texdoc.net/texmf-dist/doc/latex/base/fixltx2e.pdf
\documentclass[ 
  twoside,openright,titlepage,numbers=noenddot,headinclude,%1headlines,% 
  %%letterpaper a4paper
  footinclude=true,cleardoublepage=empty,abstractoff, % <--- obsolete, remove (todo)
  BCOR=5mm,paper=letter,fontsize=11pt,
  ngerman,american,
  ]{scrreprt}

% classicthesis-config.tex 
% Use it at the beginning of your ClassicThesis.tex, or as a LaTeX Preamble 
% in your ClassicThesis.{tex,lyx} with % classicthesis-config.tex 
% Use it at the beginning of your ClassicThesis.tex, or as a LaTeX Preamble 
% in your ClassicThesis.{tex,lyx} with % classicthesis-config.tex 
% Use it at the beginning of your ClassicThesis.tex, or as a LaTeX Preamble 
% in your ClassicThesis.{tex,lyx} with \input{classicthesis-config}

\PassOptionsToPackage{utf8}{inputenc}
\usepackage{inputenc}

% *****************************************************************************
% 1. Configure classicthesis for your needs here, e.g., remove "drafting" 
% below in order to deactivate the time-stamp on the pages
% *****************************************************************************
\PassOptionsToPackage{eulerchapternumbers,listings,dottedtoc,
  pdfspacing,%floatperchapter,linedheaders,drafting
  subfig,beramono,parts}{classicthesis}

% ********************************************************************
% Available options for classicthesis.sty 
% (see ClassicThesis.pdf for more information):
% drafting
% parts nochapters linedheaders
% eulerchapternumbers beramono eulermath pdfspacing minionprospacing
% tocaligned dottedtoc manychapters
% listings floatperchapter subfig
% ********************************************************************

% *****************************************************************************
% 2. Personal data and user ad-hoc commands
% *****************************************************************************
\newcommand{\myTitle}{Molecular Motion Movies and Geometry Reconstruction Using Bayesian Statistics\xspace}
\newcommand{\mySubtitle}{Hakuna matata?\xspace}
\newcommand{\myDegree}{Master of Science (M.Sc.)\xspace}
\newcommand{\myName}{Ali Ramadhan\xspace}
\newcommand{\myProf}{Joseph Sanderson\xspace}
\newcommand{\myFaculty}{Faculty of Science\xspace}
\newcommand{\myDepartment}{Department of Physics and Astronomy\xspace}
\newcommand{\myUni}{University of Waterloo\xspace}
\newcommand{\myLocation}{Waterloo, Ontario, Canada\xspace}
\newcommand{\myTime}{\today\xspace}

% Setup, finetuning, and useful commands
\newcounter{dummy} % necessary for correct hyperlinks (to index, bib, etc.)
\newlength{\abcd} % for ab..z string length calculation
\providecommand{\mLyX}{L\kern-.1667em\lower.25em\hbox{Y}\kern-.125emX\@}

\newcommand{\ie}{i.\,e.}
\newcommand{\Ie}{I.\,e.}
\newcommand{\eg}{e.\,g.}
\newcommand{\Eg}{E.\,g.} 

% *****************************************************************************
% 3. Loading some handy packages
% *****************************************************************************

% Packages with options that might require adjustments
\usepackage{babel}                  

\usepackage{csquotes}
\PassOptionsToPackage{
  %backend=biber, %instead of bibtex
	backend=bibtex8,bibencoding=ascii,
	language=auto,
	style=authoryear-comp,
  bibstyle=authoryear,
  firstinits=true,
  %dashed=false, % dashed: substitute rep. author with ---
  sorting=nyt,
  maxcitenames=3,
  maxbibnames=10,
  %backref=true,
  natbib=true % natbib compatibility mode (\citep and \citet still work)
}{biblatex}

\usepackage{biblatex}

%\PassOptionsToPackage{fleqn}{amsmath} % math environments and more by the AMS 
\usepackage{amsmath}

% General useful packages
\PassOptionsToPackage{T1}{fontenc} % T2A for cyrillics
    \usepackage{fontenc}     
\usepackage{textcomp} % fix warning with missing font shapes
\usepackage{scrhack}  % fix warnings when using KOMA with listings package
\usepackage{xspace}   % to get the spacing after macros right  
\usepackage{mparhack} % get marginpar right
\usepackage{fixltx2e} % fixes some LaTeX stuff --> since 2015 in the LaTeX kernel (see below)
%\usepackage[latest]{latexrelease} % will be used once available in more distributions (ISSUE #107)

\PassOptionsToPackage{printonlyused,smaller}{acronym} 
    \usepackage{acronym} % nice macros for handling all acronyms in the thesis
    %\renewcommand{\bflabel}[1]{{#1}\hfill} % fix the list of acronyms --> no longer working
    %\renewcommand*{\acsfont}[1]{\textsc{#1}} 
    \renewcommand*{\aclabelfont}[1]{\acsfont{#1}}

% *****************************************************************************
% 4. Setup floats: tables, (sub)figures, and captions
% *****************************************************************************
\usepackage{tabularx} % better tables
    \setlength{\extrarowheight}{3pt} % increase table row height
\newcommand{\tableheadline}[1]{\multicolumn{1}{c}{\spacedlowsmallcaps{#1}}}
\newcommand{\myfloatalign}{\centering} % to be used with each float for alignment

\usepackage{caption}
%\captionsetup{font=small,labelformat=smallcaps} % format=hang,
\captionsetup{font=small, format=plain} % format=hang,

% Thanks to cgnieder and Claus Lahiri
% http://tex.stackexchange.com/questions/69349/spacedlowsmallcaps-in-caption-label
% [REMOVED DUE TO OTHER PROBLEMS, SEE ISSUE #82]    
%\DeclareCaptionLabelFormat{smallcaps}{\bothIfFirst{#1}{~}\MakeTextLowercase{\textsc{#2}}}

\usepackage{subfig}

% *****************************************************************************
% 5. Setup code listings
% *****************************************************************************
\usepackage{listings} 
%\lstset{emph={trueIndex,root},emphstyle=\color{BlueViolet}}%\underbar} % for special keywords
\lstset{language=[LaTeX]Tex,%C++,
    morekeywords={PassOptionsToPackage,selectlanguage},
    keywordstyle=\color{RoyalBlue},%\bfseries,
    basicstyle=\small\ttfamily,
    %identifierstyle=\color{NavyBlue},
    commentstyle=\color{Green}\ttfamily,
    stringstyle=\rmfamily,
    numbers=none,%left,%
    numberstyle=\scriptsize,%\tiny
    stepnumber=5,
    numbersep=8pt,
    showstringspaces=false,
    breaklines=true,
    %frameround=ftff,
    %frame=single,
    belowcaptionskip=.75\baselineskip
    %frame=L
} 

% *****************************************************************************
% 6. PDFLaTeX, hyperreferences and citation backreferences
% *****************************************************************************

% Using PDFLaTeX
\PassOptionsToPackage{pdftex,hyperfootnotes=false,pdfpagelabels}{hyperref}
    \usepackage{hyperref}  % backref linktocpage pagebackref
\pdfcompresslevel=9
\pdfadjustspacing=1 
\PassOptionsToPackage{pdftex}{graphicx}
    \usepackage{graphicx}

% Hyperreferences
\hypersetup{%
    %draft, % = no hyperlinking at all (useful in b/w printouts)
    colorlinks=true, linktocpage=true, pdfstartpage=3, pdfstartview=FitV,%
    % uncomment the following line if you want to have black links (e.g., for printing)
    %colorlinks=false, linktocpage=false, pdfstartpage=3, pdfstartview=FitV, pdfborder={0 0 0},%
    breaklinks=true, pdfpagemode=UseNone, pageanchor=true, pdfpagemode=UseOutlines,%
    plainpages=false, bookmarksnumbered, bookmarksopen=true, bookmarksopenlevel=1,%
    hypertexnames=true, pdfhighlight=/O,%nesting=true,%frenchlinks,%
    urlcolor=webbrown, linkcolor=RoyalBlue, citecolor=webgreen, %pagecolor=RoyalBlue,%
    %urlcolor=Black, linkcolor=Black, citecolor=Black, %pagecolor=Black,%
    pdftitle={\myTitle},%
    pdfauthor={\textcopyright\ \myName, \myUni, \myFaculty},%
    pdfsubject={},%
    pdfkeywords={},%
    pdfcreator={pdfLaTeX},%
    pdfproducer={LaTeX with hyperref and classicthesis}%
}   

% Setup autoreferences
% There are some issues regarding autorefnames
% http://www.ureader.de/msg/136221647.aspx
% http://www.tex.ac.uk/cgi-bin/texfaq2html?label=latexwords
% you have to redefine the makros for the 
% language you use, e.g., american, ngerman
% (as chosen when loading babel/AtBeginDocument)
\makeatletter
\@ifpackageloaded{babel}%
    {%
       \addto\extrasamerican{%
			\renewcommand*{\figureautorefname}{Figure}%
			\renewcommand*{\tableautorefname}{Table}%
			\renewcommand*{\partautorefname}{Part}%
			\renewcommand*{\chapterautorefname}{Chapter}%
			\renewcommand*{\sectionautorefname}{Section}%
			\renewcommand*{\subsectionautorefname}{Section}%
			\renewcommand*{\subsubsectionautorefname}{Section}%     
                }%
       \addto\extrasngerman{% 
			\renewcommand*{\paragraphautorefname}{Absatz}%
			\renewcommand*{\subparagraphautorefname}{Unterabsatz}%
			\renewcommand*{\footnoteautorefname}{Fu\"snote}%
			\renewcommand*{\FancyVerbLineautorefname}{Zeile}%
			\renewcommand*{\theoremautorefname}{Theorem}%
			\renewcommand*{\appendixautorefname}{Anhang}%
			\renewcommand*{\equationautorefname}{Gleichung}%        
			\renewcommand*{\itemautorefname}{Punkt}%
                }%  
            % Fix to getting autorefs for subfigures right (thanks to Belinda Vogt for changing the definition)
            \providecommand{\subfigureautorefname}{\figureautorefname}     
    }{\relax}
\makeatother

% *****************************************************************************
% 7. Last calls before the bar closes
% *****************************************************************************

% Development Stuff
\listfiles
%\PassOptionsToPackage{l2tabu,orthodox,abort}{nag}
%   \usepackage{nag}
%\PassOptionsToPackage{warning, all}{onlyamsmath}
%   \usepackage{onlyamsmath}

% Last, but not least...
\usepackage{classicthesis} 

% *****************************************************************************
% 8. Further adjustments (experimental)
% *****************************************************************************

% Changing the text area
%\linespread{1.05} % a bit more for Palatino
%\areaset[current]{312pt}{761pt} % 686 (factor 2.2) + 33 head + 42 head \the\footskip
%\setlength{\marginparwidth}{7em}%
%\setlength{\marginparsep}{2em}%

% Using different fonts
%\usepackage[oldstylenums]{kpfonts} % oldstyle notextcomp
%\usepackage[osf]{libertine}
%\usepackage[light,condensed,math]{iwona}
%\renewcommand{\sfdefault}{iwona}
%\usepackage{lmodern} % <-- no osf support :-(
%\usepackage{cfr-lm} % 
%\usepackage[urw-garamond]{mathdesign} <-- no osf support :-(
%\usepackage[default,osfigures]{opensans} % scale=0.95 
%\usepackage[sfdefault]{FiraSans}

% *****************************************************************************
% 9. My own extra configuration.
% *****************************************************************************

% Use \sum symbol from Computer Modern (LaTeX default) math font.
\DeclareSymbolFont{cmlargesymbols}{OMX}{cmex}{m}{n}
\let\sum\relax
\DeclareMathSymbol{\sum}{\mathop}{cmlargesymbols}{"50}

\numberwithin{equation}{chapter}

% Theorems and definitions
\usepackage{amsthm}
\newtheorem{theorem}{Theorem}[chapter]
\newtheorem{definition}{Definition}[chapter]

%\usepackage[tight]{minitoc}
\usepackage{minitoc}
\nomtcrule
%\renewcommand{\mtifont}{\normalsize\textsc}

\renewbibmacro{in:}{}
%\renewcommand{\mkbibnamefamily}[1]{\textsc{#1}} % Small caps for author names.

\usepackage{imakeidx}
\usepackage{sidecap}
%\usepackage[version=4]{mhchem}
\usepackage{chemformula}
\usepackage{float} % For [H] option in figures.
\usepackage{gensymb} % For \degree
%\usepackage[toc,page]{appendix}
\usepackage{setspace}

% to automatically ensure floats do not go into the next section.
\usepackage[section]{placeins}

% Unit typesetting
\usepackage[binary-units=true]{siunitx}
\sisetup{list-final-separator = {, and }}

\PassOptionsToPackage{utf8}{inputenc}
\usepackage{inputenc}

% *****************************************************************************
% 1. Configure classicthesis for your needs here, e.g., remove "drafting" 
% below in order to deactivate the time-stamp on the pages
% *****************************************************************************
\PassOptionsToPackage{eulerchapternumbers,listings,dottedtoc,
  pdfspacing,%floatperchapter,linedheaders,drafting
  subfig,beramono,parts}{classicthesis}

% ********************************************************************
% Available options for classicthesis.sty 
% (see ClassicThesis.pdf for more information):
% drafting
% parts nochapters linedheaders
% eulerchapternumbers beramono eulermath pdfspacing minionprospacing
% tocaligned dottedtoc manychapters
% listings floatperchapter subfig
% ********************************************************************

% *****************************************************************************
% 2. Personal data and user ad-hoc commands
% *****************************************************************************
\newcommand{\myTitle}{Molecular Motion Movies and Geometry Reconstruction Using Bayesian Statistics\xspace}
\newcommand{\mySubtitle}{Hakuna matata?\xspace}
\newcommand{\myDegree}{Master of Science (M.Sc.)\xspace}
\newcommand{\myName}{Ali Ramadhan\xspace}
\newcommand{\myProf}{Joseph Sanderson\xspace}
\newcommand{\myFaculty}{Faculty of Science\xspace}
\newcommand{\myDepartment}{Department of Physics and Astronomy\xspace}
\newcommand{\myUni}{University of Waterloo\xspace}
\newcommand{\myLocation}{Waterloo, Ontario, Canada\xspace}
\newcommand{\myTime}{\today\xspace}

% Setup, finetuning, and useful commands
\newcounter{dummy} % necessary for correct hyperlinks (to index, bib, etc.)
\newlength{\abcd} % for ab..z string length calculation
\providecommand{\mLyX}{L\kern-.1667em\lower.25em\hbox{Y}\kern-.125emX\@}

\newcommand{\ie}{i.\,e.}
\newcommand{\Ie}{I.\,e.}
\newcommand{\eg}{e.\,g.}
\newcommand{\Eg}{E.\,g.} 

% *****************************************************************************
% 3. Loading some handy packages
% *****************************************************************************

% Packages with options that might require adjustments
\usepackage{babel}                  

\usepackage{csquotes}
\PassOptionsToPackage{
  %backend=biber, %instead of bibtex
	backend=bibtex8,bibencoding=ascii,
	language=auto,
	style=authoryear-comp,
  bibstyle=authoryear,
  firstinits=true,
  %dashed=false, % dashed: substitute rep. author with ---
  sorting=nyt,
  maxcitenames=3,
  maxbibnames=10,
  %backref=true,
  natbib=true % natbib compatibility mode (\citep and \citet still work)
}{biblatex}

\usepackage{biblatex}

%\PassOptionsToPackage{fleqn}{amsmath} % math environments and more by the AMS 
\usepackage{amsmath}

% General useful packages
\PassOptionsToPackage{T1}{fontenc} % T2A for cyrillics
    \usepackage{fontenc}     
\usepackage{textcomp} % fix warning with missing font shapes
\usepackage{scrhack}  % fix warnings when using KOMA with listings package
\usepackage{xspace}   % to get the spacing after macros right  
\usepackage{mparhack} % get marginpar right
\usepackage{fixltx2e} % fixes some LaTeX stuff --> since 2015 in the LaTeX kernel (see below)
%\usepackage[latest]{latexrelease} % will be used once available in more distributions (ISSUE #107)

\PassOptionsToPackage{printonlyused,smaller}{acronym} 
    \usepackage{acronym} % nice macros for handling all acronyms in the thesis
    %\renewcommand{\bflabel}[1]{{#1}\hfill} % fix the list of acronyms --> no longer working
    %\renewcommand*{\acsfont}[1]{\textsc{#1}} 
    \renewcommand*{\aclabelfont}[1]{\acsfont{#1}}

% *****************************************************************************
% 4. Setup floats: tables, (sub)figures, and captions
% *****************************************************************************
\usepackage{tabularx} % better tables
    \setlength{\extrarowheight}{3pt} % increase table row height
\newcommand{\tableheadline}[1]{\multicolumn{1}{c}{\spacedlowsmallcaps{#1}}}
\newcommand{\myfloatalign}{\centering} % to be used with each float for alignment

\usepackage{caption}
%\captionsetup{font=small,labelformat=smallcaps} % format=hang,
\captionsetup{font=small, format=plain} % format=hang,

% Thanks to cgnieder and Claus Lahiri
% http://tex.stackexchange.com/questions/69349/spacedlowsmallcaps-in-caption-label
% [REMOVED DUE TO OTHER PROBLEMS, SEE ISSUE #82]    
%\DeclareCaptionLabelFormat{smallcaps}{\bothIfFirst{#1}{~}\MakeTextLowercase{\textsc{#2}}}

\usepackage{subfig}

% *****************************************************************************
% 5. Setup code listings
% *****************************************************************************
\usepackage{listings} 
%\lstset{emph={trueIndex,root},emphstyle=\color{BlueViolet}}%\underbar} % for special keywords
\lstset{language=[LaTeX]Tex,%C++,
    morekeywords={PassOptionsToPackage,selectlanguage},
    keywordstyle=\color{RoyalBlue},%\bfseries,
    basicstyle=\small\ttfamily,
    %identifierstyle=\color{NavyBlue},
    commentstyle=\color{Green}\ttfamily,
    stringstyle=\rmfamily,
    numbers=none,%left,%
    numberstyle=\scriptsize,%\tiny
    stepnumber=5,
    numbersep=8pt,
    showstringspaces=false,
    breaklines=true,
    %frameround=ftff,
    %frame=single,
    belowcaptionskip=.75\baselineskip
    %frame=L
} 

% *****************************************************************************
% 6. PDFLaTeX, hyperreferences and citation backreferences
% *****************************************************************************

% Using PDFLaTeX
\PassOptionsToPackage{pdftex,hyperfootnotes=false,pdfpagelabels}{hyperref}
    \usepackage{hyperref}  % backref linktocpage pagebackref
\pdfcompresslevel=9
\pdfadjustspacing=1 
\PassOptionsToPackage{pdftex}{graphicx}
    \usepackage{graphicx}

% Hyperreferences
\hypersetup{%
    %draft, % = no hyperlinking at all (useful in b/w printouts)
    colorlinks=true, linktocpage=true, pdfstartpage=3, pdfstartview=FitV,%
    % uncomment the following line if you want to have black links (e.g., for printing)
    %colorlinks=false, linktocpage=false, pdfstartpage=3, pdfstartview=FitV, pdfborder={0 0 0},%
    breaklinks=true, pdfpagemode=UseNone, pageanchor=true, pdfpagemode=UseOutlines,%
    plainpages=false, bookmarksnumbered, bookmarksopen=true, bookmarksopenlevel=1,%
    hypertexnames=true, pdfhighlight=/O,%nesting=true,%frenchlinks,%
    urlcolor=webbrown, linkcolor=RoyalBlue, citecolor=webgreen, %pagecolor=RoyalBlue,%
    %urlcolor=Black, linkcolor=Black, citecolor=Black, %pagecolor=Black,%
    pdftitle={\myTitle},%
    pdfauthor={\textcopyright\ \myName, \myUni, \myFaculty},%
    pdfsubject={},%
    pdfkeywords={},%
    pdfcreator={pdfLaTeX},%
    pdfproducer={LaTeX with hyperref and classicthesis}%
}   

% Setup autoreferences
% There are some issues regarding autorefnames
% http://www.ureader.de/msg/136221647.aspx
% http://www.tex.ac.uk/cgi-bin/texfaq2html?label=latexwords
% you have to redefine the makros for the 
% language you use, e.g., american, ngerman
% (as chosen when loading babel/AtBeginDocument)
\makeatletter
\@ifpackageloaded{babel}%
    {%
       \addto\extrasamerican{%
			\renewcommand*{\figureautorefname}{Figure}%
			\renewcommand*{\tableautorefname}{Table}%
			\renewcommand*{\partautorefname}{Part}%
			\renewcommand*{\chapterautorefname}{Chapter}%
			\renewcommand*{\sectionautorefname}{Section}%
			\renewcommand*{\subsectionautorefname}{Section}%
			\renewcommand*{\subsubsectionautorefname}{Section}%     
                }%
       \addto\extrasngerman{% 
			\renewcommand*{\paragraphautorefname}{Absatz}%
			\renewcommand*{\subparagraphautorefname}{Unterabsatz}%
			\renewcommand*{\footnoteautorefname}{Fu\"snote}%
			\renewcommand*{\FancyVerbLineautorefname}{Zeile}%
			\renewcommand*{\theoremautorefname}{Theorem}%
			\renewcommand*{\appendixautorefname}{Anhang}%
			\renewcommand*{\equationautorefname}{Gleichung}%        
			\renewcommand*{\itemautorefname}{Punkt}%
                }%  
            % Fix to getting autorefs for subfigures right (thanks to Belinda Vogt for changing the definition)
            \providecommand{\subfigureautorefname}{\figureautorefname}     
    }{\relax}
\makeatother

% *****************************************************************************
% 7. Last calls before the bar closes
% *****************************************************************************

% Development Stuff
\listfiles
%\PassOptionsToPackage{l2tabu,orthodox,abort}{nag}
%   \usepackage{nag}
%\PassOptionsToPackage{warning, all}{onlyamsmath}
%   \usepackage{onlyamsmath}

% Last, but not least...
\usepackage{classicthesis} 

% *****************************************************************************
% 8. Further adjustments (experimental)
% *****************************************************************************

% Changing the text area
%\linespread{1.05} % a bit more for Palatino
%\areaset[current]{312pt}{761pt} % 686 (factor 2.2) + 33 head + 42 head \the\footskip
%\setlength{\marginparwidth}{7em}%
%\setlength{\marginparsep}{2em}%

% Using different fonts
%\usepackage[oldstylenums]{kpfonts} % oldstyle notextcomp
%\usepackage[osf]{libertine}
%\usepackage[light,condensed,math]{iwona}
%\renewcommand{\sfdefault}{iwona}
%\usepackage{lmodern} % <-- no osf support :-(
%\usepackage{cfr-lm} % 
%\usepackage[urw-garamond]{mathdesign} <-- no osf support :-(
%\usepackage[default,osfigures]{opensans} % scale=0.95 
%\usepackage[sfdefault]{FiraSans}

% *****************************************************************************
% 9. My own extra configuration.
% *****************************************************************************

% Use \sum symbol from Computer Modern (LaTeX default) math font.
\DeclareSymbolFont{cmlargesymbols}{OMX}{cmex}{m}{n}
\let\sum\relax
\DeclareMathSymbol{\sum}{\mathop}{cmlargesymbols}{"50}

\numberwithin{equation}{chapter}

% Theorems and definitions
\usepackage{amsthm}
\newtheorem{theorem}{Theorem}[chapter]
\newtheorem{definition}{Definition}[chapter]

%\usepackage[tight]{minitoc}
\usepackage{minitoc}
\nomtcrule
%\renewcommand{\mtifont}{\normalsize\textsc}

\renewbibmacro{in:}{}
%\renewcommand{\mkbibnamefamily}[1]{\textsc{#1}} % Small caps for author names.

\usepackage{imakeidx}
\usepackage{sidecap}
%\usepackage[version=4]{mhchem}
\usepackage{chemformula}
\usepackage{float} % For [H] option in figures.
\usepackage{gensymb} % For \degree
%\usepackage[toc,page]{appendix}
\usepackage{setspace}

% to automatically ensure floats do not go into the next section.
\usepackage[section]{placeins}

% Unit typesetting
\usepackage[binary-units=true]{siunitx}
\sisetup{list-final-separator = {, and }}

\PassOptionsToPackage{utf8}{inputenc}
\usepackage{inputenc}

% *****************************************************************************
% 1. Configure classicthesis for your needs here, e.g., remove "drafting" 
% below in order to deactivate the time-stamp on the pages
% *****************************************************************************
\PassOptionsToPackage{eulerchapternumbers,listings,dottedtoc,
  pdfspacing,%floatperchapter,linedheaders,drafting
  subfig,beramono,parts}{classicthesis}

% ********************************************************************
% Available options for classicthesis.sty 
% (see ClassicThesis.pdf for more information):
% drafting
% parts nochapters linedheaders
% eulerchapternumbers beramono eulermath pdfspacing minionprospacing
% tocaligned dottedtoc manychapters
% listings floatperchapter subfig
% ********************************************************************

% *****************************************************************************
% 2. Personal data and user ad-hoc commands
% *****************************************************************************
\newcommand{\myTitle}{Molecular Motion Movies and Geometry Reconstruction Using Bayesian Statistics\xspace}
\newcommand{\mySubtitle}{Hakuna matata?\xspace}
\newcommand{\myDegree}{Master of Science (M.Sc.)\xspace}
\newcommand{\myName}{Ali Ramadhan\xspace}
\newcommand{\myProf}{Joseph Sanderson\xspace}
\newcommand{\myFaculty}{Faculty of Science\xspace}
\newcommand{\myDepartment}{Department of Physics and Astronomy\xspace}
\newcommand{\myUni}{University of Waterloo\xspace}
\newcommand{\myLocation}{Waterloo, Ontario, Canada\xspace}
\newcommand{\myTime}{\today\xspace}

% Setup, finetuning, and useful commands
\newcounter{dummy} % necessary for correct hyperlinks (to index, bib, etc.)
\newlength{\abcd} % for ab..z string length calculation
\providecommand{\mLyX}{L\kern-.1667em\lower.25em\hbox{Y}\kern-.125emX\@}

\newcommand{\ie}{i.\,e.}
\newcommand{\Ie}{I.\,e.}
\newcommand{\eg}{e.\,g.}
\newcommand{\Eg}{E.\,g.} 

% *****************************************************************************
% 3. Loading some handy packages
% *****************************************************************************

% Packages with options that might require adjustments
\usepackage{babel}                  

\usepackage{csquotes}
\PassOptionsToPackage{
  %backend=biber, %instead of bibtex
	backend=bibtex8,bibencoding=ascii,
	language=auto,
	style=authoryear-comp,
  bibstyle=authoryear,
  firstinits=true,
  %dashed=false, % dashed: substitute rep. author with ---
  sorting=nyt,
  maxcitenames=3,
  maxbibnames=10,
  %backref=true,
  natbib=true % natbib compatibility mode (\citep and \citet still work)
}{biblatex}

\usepackage{biblatex}

%\PassOptionsToPackage{fleqn}{amsmath} % math environments and more by the AMS 
\usepackage{amsmath}

% General useful packages
\PassOptionsToPackage{T1}{fontenc} % T2A for cyrillics
    \usepackage{fontenc}     
\usepackage{textcomp} % fix warning with missing font shapes
\usepackage{scrhack}  % fix warnings when using KOMA with listings package
\usepackage{xspace}   % to get the spacing after macros right  
\usepackage{mparhack} % get marginpar right
\usepackage{fixltx2e} % fixes some LaTeX stuff --> since 2015 in the LaTeX kernel (see below)
%\usepackage[latest]{latexrelease} % will be used once available in more distributions (ISSUE #107)

\PassOptionsToPackage{printonlyused,smaller}{acronym} 
    \usepackage{acronym} % nice macros for handling all acronyms in the thesis
    %\renewcommand{\bflabel}[1]{{#1}\hfill} % fix the list of acronyms --> no longer working
    %\renewcommand*{\acsfont}[1]{\textsc{#1}} 
    \renewcommand*{\aclabelfont}[1]{\acsfont{#1}}

% *****************************************************************************
% 4. Setup floats: tables, (sub)figures, and captions
% *****************************************************************************
\usepackage{tabularx} % better tables
    \setlength{\extrarowheight}{3pt} % increase table row height
\newcommand{\tableheadline}[1]{\multicolumn{1}{c}{\spacedlowsmallcaps{#1}}}
\newcommand{\myfloatalign}{\centering} % to be used with each float for alignment

\usepackage{caption}
%\captionsetup{font=small,labelformat=smallcaps} % format=hang,
\captionsetup{font=small, format=plain} % format=hang,

% Thanks to cgnieder and Claus Lahiri
% http://tex.stackexchange.com/questions/69349/spacedlowsmallcaps-in-caption-label
% [REMOVED DUE TO OTHER PROBLEMS, SEE ISSUE #82]    
%\DeclareCaptionLabelFormat{smallcaps}{\bothIfFirst{#1}{~}\MakeTextLowercase{\textsc{#2}}}

\usepackage{subfig}

% *****************************************************************************
% 5. Setup code listings
% *****************************************************************************
\usepackage{listings} 
%\lstset{emph={trueIndex,root},emphstyle=\color{BlueViolet}}%\underbar} % for special keywords
\lstset{language=[LaTeX]Tex,%C++,
    morekeywords={PassOptionsToPackage,selectlanguage},
    keywordstyle=\color{RoyalBlue},%\bfseries,
    basicstyle=\small\ttfamily,
    %identifierstyle=\color{NavyBlue},
    commentstyle=\color{Green}\ttfamily,
    stringstyle=\rmfamily,
    numbers=none,%left,%
    numberstyle=\scriptsize,%\tiny
    stepnumber=5,
    numbersep=8pt,
    showstringspaces=false,
    breaklines=true,
    %frameround=ftff,
    %frame=single,
    belowcaptionskip=.75\baselineskip
    %frame=L
} 

% *****************************************************************************
% 6. PDFLaTeX, hyperreferences and citation backreferences
% *****************************************************************************

% Using PDFLaTeX
\PassOptionsToPackage{pdftex,hyperfootnotes=false,pdfpagelabels}{hyperref}
    \usepackage{hyperref}  % backref linktocpage pagebackref
\pdfcompresslevel=9
\pdfadjustspacing=1 
\PassOptionsToPackage{pdftex}{graphicx}
    \usepackage{graphicx}

% Hyperreferences
\hypersetup{%
    %draft, % = no hyperlinking at all (useful in b/w printouts)
    colorlinks=true, linktocpage=true, pdfstartpage=3, pdfstartview=FitV,%
    % uncomment the following line if you want to have black links (e.g., for printing)
    %colorlinks=false, linktocpage=false, pdfstartpage=3, pdfstartview=FitV, pdfborder={0 0 0},%
    breaklinks=true, pdfpagemode=UseNone, pageanchor=true, pdfpagemode=UseOutlines,%
    plainpages=false, bookmarksnumbered, bookmarksopen=true, bookmarksopenlevel=1,%
    hypertexnames=true, pdfhighlight=/O,%nesting=true,%frenchlinks,%
    urlcolor=webbrown, linkcolor=RoyalBlue, citecolor=webgreen, %pagecolor=RoyalBlue,%
    %urlcolor=Black, linkcolor=Black, citecolor=Black, %pagecolor=Black,%
    pdftitle={\myTitle},%
    pdfauthor={\textcopyright\ \myName, \myUni, \myFaculty},%
    pdfsubject={},%
    pdfkeywords={},%
    pdfcreator={pdfLaTeX},%
    pdfproducer={LaTeX with hyperref and classicthesis}%
}   

% Setup autoreferences
% There are some issues regarding autorefnames
% http://www.ureader.de/msg/136221647.aspx
% http://www.tex.ac.uk/cgi-bin/texfaq2html?label=latexwords
% you have to redefine the makros for the 
% language you use, e.g., american, ngerman
% (as chosen when loading babel/AtBeginDocument)
\makeatletter
\@ifpackageloaded{babel}%
    {%
       \addto\extrasamerican{%
			\renewcommand*{\figureautorefname}{Figure}%
			\renewcommand*{\tableautorefname}{Table}%
			\renewcommand*{\partautorefname}{Part}%
			\renewcommand*{\chapterautorefname}{Chapter}%
			\renewcommand*{\sectionautorefname}{Section}%
			\renewcommand*{\subsectionautorefname}{Section}%
			\renewcommand*{\subsubsectionautorefname}{Section}%     
                }%
       \addto\extrasngerman{% 
			\renewcommand*{\paragraphautorefname}{Absatz}%
			\renewcommand*{\subparagraphautorefname}{Unterabsatz}%
			\renewcommand*{\footnoteautorefname}{Fu\"snote}%
			\renewcommand*{\FancyVerbLineautorefname}{Zeile}%
			\renewcommand*{\theoremautorefname}{Theorem}%
			\renewcommand*{\appendixautorefname}{Anhang}%
			\renewcommand*{\equationautorefname}{Gleichung}%        
			\renewcommand*{\itemautorefname}{Punkt}%
                }%  
            % Fix to getting autorefs for subfigures right (thanks to Belinda Vogt for changing the definition)
            \providecommand{\subfigureautorefname}{\figureautorefname}     
    }{\relax}
\makeatother

% *****************************************************************************
% 7. Last calls before the bar closes
% *****************************************************************************

% Development Stuff
\listfiles
%\PassOptionsToPackage{l2tabu,orthodox,abort}{nag}
%   \usepackage{nag}
%\PassOptionsToPackage{warning, all}{onlyamsmath}
%   \usepackage{onlyamsmath}

% Last, but not least...
\usepackage{classicthesis} 

% *****************************************************************************
% 8. Further adjustments (experimental)
% *****************************************************************************

% Changing the text area
%\linespread{1.05} % a bit more for Palatino
%\areaset[current]{312pt}{761pt} % 686 (factor 2.2) + 33 head + 42 head \the\footskip
%\setlength{\marginparwidth}{7em}%
%\setlength{\marginparsep}{2em}%

% Using different fonts
%\usepackage[oldstylenums]{kpfonts} % oldstyle notextcomp
%\usepackage[osf]{libertine}
%\usepackage[light,condensed,math]{iwona}
%\renewcommand{\sfdefault}{iwona}
%\usepackage{lmodern} % <-- no osf support :-(
%\usepackage{cfr-lm} % 
%\usepackage[urw-garamond]{mathdesign} <-- no osf support :-(
%\usepackage[default,osfigures]{opensans} % scale=0.95 
%\usepackage[sfdefault]{FiraSans}

% *****************************************************************************
% 9. My own extra configuration.
% *****************************************************************************

% Use \sum symbol from Computer Modern (LaTeX default) math font.
\DeclareSymbolFont{cmlargesymbols}{OMX}{cmex}{m}{n}
\let\sum\relax
\DeclareMathSymbol{\sum}{\mathop}{cmlargesymbols}{"50}

\numberwithin{equation}{chapter}

% Theorems and definitions
\usepackage{amsthm}
\newtheorem{theorem}{Theorem}[chapter]
\newtheorem{definition}{Definition}[chapter]

%\usepackage[tight]{minitoc}
\usepackage{minitoc}
\nomtcrule
%\renewcommand{\mtifont}{\normalsize\textsc}

\renewbibmacro{in:}{}
%\renewcommand{\mkbibnamefamily}[1]{\textsc{#1}} % Small caps for author names.

\usepackage{imakeidx}
\usepackage{sidecap}
%\usepackage[version=4]{mhchem}
\usepackage{chemformula}
\usepackage{float} % For [H] option in figures.
\usepackage{gensymb} % For \degree
%\usepackage[toc,page]{appendix}
\usepackage{setspace}

% to automatically ensure floats do not go into the next section.
\usepackage[section]{placeins}

% Unit typesetting
\usepackage[binary-units=true]{siunitx}
\sisetup{list-final-separator = {, and }}

\numberwithin{equation}{chapter}

\addbibresource{Bibliographies/CEI.bib}
\addbibresource{Bibliographies/MolecularMovies.bib}

%\hyphenation{put special hyphenation here}

\begin{document}
\frenchspacing
\raggedbottom
\selectlanguage{american} % american ngerman
%\renewcommand*{\bibname}{new name}
%\setbibpreamble{}
\pagenumbering{roman}
\pagestyle{plain}

%*******************************************************
% Little Dirty Titlepage
%*******************************************************
\thispagestyle{empty}
%\pdfbookmark[1]{Titel}{title}
%*******************************************************
\begin{center}
    \spacedlowsmallcaps{\myName} \\ \medskip                        

    \begingroup
        \color{Maroon}\spacedallcaps{\myTitle}
    \endgroup
\end{center}        

%*******************************************************
% Titlepage
%*******************************************************
\begin{titlepage}
    % if you want the titlepage to be centered, uncomment and fine-tune the line below (KOMA classes environment)
    \begin{addmargin}[-1cm]{-3cm}
    \begin{center}
        \large  

        \hfill

        \vfill

        \begingroup
            \color{Maroon}\spacedallcaps{\myTitle} \\ \bigskip
        \endgroup

        \spacedlowsmallcaps{\myName}

        \vfill

        \includegraphics[width=16cm]{gfx/simba} \\ \medskip

        \mySubtitle \\ \medskip   
        %\myDegree \\
        %\myDepartment \\                            
        %\myFaculty \\
        %\myUni \\ \bigskip

        \myTime\ %-- \myVersion

        \vfill                      

    \end{center}  
  \end{addmargin}       
\end{titlepage}
\thispagestyle{empty}

\hfill

\vfill

\noindent\myName: \textit{\myTitle.} %\mySubtitle, %\myDegree, 
\textcopyright\ \myTime

%\bigskip
%
%\noindent\spacedlowsmallcaps{Supervisors}: \\
%\myProf \\
%\myOtherProf \\ 
%\mySupervisor
%
%\medskip
%
%\noindent\spacedlowsmallcaps{Location}: \\
%\myLocation
%
%\medskip
%
%\noindent\spacedlowsmallcaps{Time Frame}: \\
%\myTime

%\cleardoublepage%*******************************************************
% Dedication
%*******************************************************
\thispagestyle{empty}
%\phantomsection 
\refstepcounter{dummy}
\pdfbookmark[1]{Dedication}{Dedication}

\vspace*{3cm}

\begin{center}
    \emph{Ohana} means family. \\
    Family means nobody gets left behind, or forgotten. \\ \medskip
    --- Lilo \& Stitch    
\end{center}

\medskip

\begin{center}
    Dedicated to the loving memory of Rudolf Miede. \\ \smallskip
    1939\,--\,2005
\end{center}
%\cleardoublepage\include{FrontBackmatter/Foreword}
\cleardoublepage%\renewcommand{\abstractname}{Abstract}
\pdfbookmark[1]{Abstract}{Abstract}
\begingroup
\let\clearpage\relax
\let\cleardoublepage\relax
\let\cleardoublepage\relax

\chapter*{Technical abstract}
Ever since the early days of Coulomb explosion imaging, the direct imaging of the structure and dynamics of individual molecules has been promised; however, in practice no one has been able to accurately retrieve this structure and all imaging relies on plotting the momentum vectors of the fragments in different ways to infer crude changes in geometry. The momentum vectors tell a large part of the story but do not provide a clear picture of the molecular dynamics everyone seeks so an accurate method of retrieving the structure is highly desirable. I use Bayesian inference and Markov chain Monte Carlo methods to elucidate the molecular geometries, allowing for the study of larger systems than before and for the inclusion of measurement error in these studies for the first time. Using this method, the ultrafast isomerization of carbonyl sulfide and acetylene is imaged in position-space and showcased as a molecular movie. The method, however, is much larger in scope and is generalized to any problem involving a physical system described by a system of ordinary differential equations, where the final conditions are experimentally measured and inference of the initial conditions is desirable.

\vfill

\endgroup			

\vfill
\cleardoublepage%\renewcommand{\abstractname}{Abstract}
\pdfbookmark[1]{Abstract}{Abstract}
\begingroup
\let\clearpage\relax
\let\cleardoublepage\relax
\let\cleardoublepage\relax

\chapter*{Technical abstract}
Insert better abstract.

\vfill

\endgroup			

\vfill
%\cleardoublepage%*******************************************************
% Publications
%*******************************************************
\pdfbookmark[1]{Publications}{publications}
\chapter*{Publications}\graffito{This is just an early --~and currently ugly~-- test!}
This might come in handy for PhD theses: some ideas and figures have appeared previously in the following publications:

%\noindent Put your publications from the thesis here. The packages \texttt{multibib} or \texttt{bibtopic} etc. can be used to handle multiple different bibliographies in your document.

\begin{refsection}[ownpubs]
    \small
    \nocite{*} % is local to to the enclosing refsection
    \printbibliography[heading=none]
\end{refsection}

\emph{Attention}: This requires a separate run of \texttt{bibtex} for your \texttt{refsection}, \eg, \texttt{ClassicThesis1-blx} for this file. You might also use \texttt{biber} as the backend for \texttt{biblatex}. See also \url{http://tex.stackexchange.com/questions/128196/problem-with-refsection}.
%\cleardoublepage\pdfbookmark[1]{Acknowledgments}{acknowledgments}

\begingroup
\let\clearpage\relax
\let\cleardoublepage\relax
\let\cleardoublepage\relax
\chapter*{Acknowledgments}
Hi everyone, thanks!

\endgroup
\pagestyle{scrheadings}
\cleardoublepage%*******************************************************
% Table of Contents
%*******************************************************
%\phantomsection
\refstepcounter{dummy}
\pdfbookmark[1]{\contentsname}{tableofcontents}
\setcounter{tocdepth}{2} % <-- 2 includes up to subsections in the ToC
\setcounter{secnumdepth}{3} % <-- 3 numbers up to subsubsections
\manualmark
\markboth{\spacedlowsmallcaps{\contentsname}}{\spacedlowsmallcaps{\contentsname}}
\tableofcontents 
\automark[section]{chapter}
\renewcommand{\chaptermark}[1]{\markboth{\spacedlowsmallcaps{#1}}{\spacedlowsmallcaps{#1}}}
\renewcommand{\sectionmark}[1]{\markright{\thesection\enspace\spacedlowsmallcaps{#1}}}
%*******************************************************
% List of Figures and of the Tables
%*******************************************************
\clearpage

\begingroup 
    \let\clearpage\relax
    \let\cleardoublepage\relax
    \let\cleardoublepage\relax
    %*******************************************************
    % List of Figures
    %*******************************************************    
    %\phantomsection 
    \refstepcounter{dummy}
    %\addcontentsline{toc}{chapter}{\listfigurename}
    \pdfbookmark[1]{\listfigurename}{lof}
    \listoffigures

    \vspace{8ex}

    %*******************************************************
    % List of Tables
    %*******************************************************
    %\phantomsection 
    \refstepcounter{dummy}
    %\addcontentsline{toc}{chapter}{\listtablename}
    \pdfbookmark[1]{\listtablename}{lot}
    \listoftables
        
    \vspace{8ex}
%   \newpage
    
    %*******************************************************
    % List of Listings
    %*******************************************************      
      %\phantomsection 
    \refstepcounter{dummy}
    %\addcontentsline{toc}{chapter}{\lstlistlistingname}
    \pdfbookmark[1]{\lstlistlistingname}{lol}
    \lstlistoflistings 

    \vspace{8ex}
       
    %*******************************************************
    % Acronyms
    %*******************************************************
    %\phantomsection 
    \refstepcounter{dummy}
    \pdfbookmark[1]{Acronyms}{acronyms}
    \markboth{\spacedlowsmallcaps{Acronyms}}{\spacedlowsmallcaps{Acronyms}}
    \chapter*{Acronyms}
    \begin{acronym}[UMLX]
        \acro{DRY}{Don't Repeat Yourself}
        \acro{API}{Application Programming Interface}
        \acro{UML}{Unified Modeling Language}
    \end{acronym}                     
\endgroup



\cleardoublepage\pagenumbering{arabic}
%\setcounter{page}{90}
% use \cleardoublepage here to avoid problems with pdfbookmark
\cleardoublepage

\ctparttext{To image the dynamics of a molecule by destroying it seems paradoxical. As we shall see, the molecule's structure is encoded in the atmoic shrapnel left behind after an explosion. However, to destroy one of nature's simplest creations is no easy task. Molecules are held together by strong chemical bonds. Our best line of attach is to shoot them with a short laser pulse---engulfing the molecule in an intense oscillating electric field will strip away its electrons and cause it to break up quickly. In these first few chapters, I will discuss how to create a short laser pulse and how these pulses interact with matter. Then I will introduce the technique of pump-probe Coulomb explosion imaging, which we will use to probe the dynamic structure of small molecules by studying the atomic fragments resulting from the explosion.}

\part{How to explode a molecule}
\chapter{The making of a short laser pulse}\label{ch:MakeShortPulse}

\section{Spontaneous and stimulated emission}
\section{Properties of a laser}
\section{Producing laser light}
\section{Making laser pulses}
\section{Making short laser pulses}
\section{Short pulse amplification}
\subsection{Chirped pulse amplification}
\subsection{Optical parametric amplification}
\section{Laser amplifiers}
\chapter{What does a short pulse do?}\label{ch:ShortPulseDoes}

\section{Short laser pulse properties}
\section{Interaction with single atoms}
\subsection{Multiphoton ionization}
\subsection{Tunneling ionization}
\subsection{Keldysh parameter}
\section{Interaction with single electrons}
\subsection{Ponderomotive force}
\subsection{Relativistic Thomson scattering}
\subsection{Nonlinear Compton scattering}
\section{Diatomic molecules in intense laser fields} 
\section{Small polyatomic molecules in intelse laser fields}

\chapter{Coulomb explosion imaging}\label{ch:CEI}

\minitoc% Creating an actual minitoc

\section{Outline of the experiment}

\subsection{Pump-probe Coulomb explosion imaging}
In pump-probe Coulomb explosion imaging (CEI) one ultrashort laser pulse is split into two pulses through the use an asymmetric beamsplitter. One of the pulses, the pump pulse, is usually much weaker than the other, the probe pulse. A time delay $\tau$ between the pulses is then created such that the pump pulse goes first and the probe pulse second. The job of the pump pulse will be to initiate some change in the molecule. One example could include an isomerization of the molecule. Thus the pump pulse ``pumps'' the molecule into some excited state. The job of the powerful probe pulse is to engulf the molecule in an intense enough laser field such that multiple electrons are stripped off of it. The molecule's individual atoms are left in a highly-charged state and begin to behave as individual point charges in a purely Coulombic potential. The entire process occurs in the presence of a constant electric field and so the positively-charged ions all accelerate upwards towards a time- and position-sensitive detector. Thus the probe pulse allows for the ``probing'' of the excited state.

\subsection{Femtosecond Multiple Pulse Length Spectroscopy}

\section{Experimental apparatus}

\section{Data analysis}
\subsection{Time and position measurement}
The position is then calculated using

\begin{equation}
x = \frac{Q_1 + Q_2}{Q_1 + Q_2 + Q_3 + Q_4} ,\quad
y = \frac{Q_1 + Q_3}{Q_1 + Q_2 + Q_3 + Q_4}
\end{equation}

\subsection{Calculating the atomic fragments' momenta}
The components of the three-dimensional momentum vector $\mathbf{p} = (p_x,p_y,p_z)$ for each atom are then calculated as

\begin{equation}
p_x = \frac{m(x-x_0)}{t} ,\;
p_y = \frac{m(y-y_0)}{t} ,\;
p_z = \frac{qV}{2\ell} \left( \frac{t_0^2 - t^2}{t} \right)
\end{equation}

\subsection{Measurement uncertainty in the momenta}
For any relation $f = f(x_1, x_2, \dots, x_n)$, assuming independent variables, the absolute uncertainty in $f$ is

\begin{equation}
df = \sqrt{\sum_{i=1}^{n} \left( \frac{\partial f}{\partial x_i} dx_i \right)^2}
\end{equation}

\section{Computationally simulating a Coulomb explosion}
To simulate an explosion of a molecule containing $n$ atoms, we must solve the classical equations of motion for each ion right after the explosion. We choose to use Hamiltonian mechanics here to acquire a system of first-order differential equations which may be easily solved by numerical methods such as the ubiquitous fourth-order Runge-Kutta. Assuming a purely electromagnetic potential for each ion, the Hamiltonian of the molecular system is
\begin{equation}
\mathcal{H}(\mathbf{r}_i, \mathbf{p}_i, t) = \sum_{i=1}^n \frac{\mathbf{p}_i^2}{2m_i} + \frac{1}{4\pi\epsilon_0}\sum_{\substack{\lbrace i,j\rbrace\\ i \ne j}} \frac{q_iq_j}{|\mathbf{r}_i-\mathbf{r}_j|}
\end{equation}
where $i,j \in \lbrace 1,2,\dots, n \rbrace$ and so the second summation is over all $i,j$ pairs where $i \ne j$. Calculating Hamilton's equations for the system, we get
\begin{subequations}
  \begin{align}
  \frac{d\mathbf{r}_i}{dt} &= \frac{\partial \mathcal{H}}{\partial \mathbf{p}_i} = \frac{\mathbf{p}_i}{m_i} \\
  \frac{d\mathbf{p}_i}{dt} &= \frac{\partial \mathcal{H}}{\partial \mathbf{r}_i} = \frac{1}{4\pi\epsilon_0}\sum_{j, \; j \ne i} \frac{\mathbf{r}_i - \mathbf{r}_j}{|\mathbf{r}_i - \mathbf{r}_j|^3}
  \end{align}
\end{subequations}
where $i$ is held fixed over the second summation.

\cleardoublepage

\ctparttext{Having destroyed a molecule and measured the momentum of each of its atomic fragments, we are left with the inverse problem of inferring its structure. While explosions proceed in a deterministic fashion, that is structures map bijectively to momentum measurements, the converse is not true. Two very different structures may produce the same momentum measurements. To make matters worse, there is no analytic solution to the problem and as the molecule grows, the problem of finding its structure becomes increasingly high dimensional. To combat this problem we will require the use of various mathematical and statistical methods. I first discuss some results from the theory of inverse problems to shed some general insight on these problems. I then follow with a discussion of optimization methods which may be used to tackle the problem for very small molecules. However, for full imaging of larger polyatomic molecules with an analysis of measurement error, Bayesian inference using Markov chain Monte Carlo methods is the way to go, which I discuss in the last chapter of this part.}

\part{How to image a molecule}
%\chapter{Bayesian inference}\label{ch:Bayesian}

\section{Introduction}
\section{Bayesian inference in the physical sciences}
\section{Probability theory}
\section{Bayes theorem}
\section{Markov chain Monte Carlo}
\section{Sampling algorithms}
\subsection{Metropolis-Hastings algorithm}
\subsection{Gibbs sampling}
\subsection{Hamiltonian MCMC} 
\section{Statistical models}
%%\addtocontents{toc}{\protect\clearpage} % <--- just debug stuff, ignore
%\include{Chapters/Chapter03}
%%\include{multiToC} % <--- just debug stuff, ignore for your documents

\ctparttext{Equipped with the means to explode a molecule and the necessary mathematical tools, we can now image the molecule's structure and shoot our molecular movie.}

\part{Shooting the movie}
\chapter{Molecular Movies}\label{ch:MolecularMovies}

\section{History and accomplishments}
The idea that Coulomb explosion imaging could be used to produce so-called molecular movies has permeated the literature ever since the technique's emergence. The goal of observing 

\subsection{Coulomb explosion imaging}

\subsubsection{Molecular movies in momentum-space}

\subsubsection{Geometry reconstruction}
\index{Classical imaging formula}
\index{Coulomb explosion imaging!Classical imaging formula}
Surprisingly, an attempt was made to arrive at an analytical solution for calculating geometries from measured momentum data. \citet{Nagaya04} \marginpar{They seem to have worked hard to find an analytical solution but their unsaid conclusion seems to be that it is an intractable problem and their group went silent on this problem.} are able to derive so-called classical imaging formulas giving the position wavefunction squared for the Coulomb explosion of a diatomic molecule and a linear triatomic molecule (the cases of symmetric and asymmetric Coulomb explosion are treated).

\index{Coulomb explosion imaging!History}
\citet{Legare05structure,Legare05dynamics} was the first to use Coulomb explosion imaging and report molecular structures. To obtain structures, they assume the explosion proceeds under a purely Coulombic potential and use optimization methods to make guesses at the structure that most accurately reproduces the observed data consistent with minimizing a least-squares objective function. Unfortunately they provide very minimal information regarding their methods and there is a complete lack of discussion acknowledging the shortcomings of this method.\marginpar{The main shortcomings being degenerate solutions and the fact that they employ convex optimization methods to a problem that is not convex. It is not clear if they even knew about these issues.} Using 8 fs laser pulses they report on the structure of D2O and SO2 \citep{Legare05structure}. They also claim to have imaged vibrating D2+ and dissociating SO22+ and SO23+ however they provide no more than a couple of dissociation frames and infer the transient D2+ bond length from kinetic energy release ratios as a function of pump-probe time delay \citep{Legare05dynamics}.

\citet{Gagnon08} reported the reconstruction of dichloromethane (CH2Cl2) using a home-made \marginpar{There is nothing wrong with writing your own code here but nonconvex optimization algorithms are tricky to get right and the reliance should be on professional code.} stochastic-based simulated annealing algorithm that globally optimizes the molecular spatial configuration. They discuss uncertainties but are only able to obtain the structure in five cases. \footnote{See \citet{Bocharova11} for a discussion of optimization algorithms.}

\index{Lookup table}
The best effort so far has probably been the one by \citet{Kunitski15} in which they use a lookup table approach to image the Efimov state of the helium trimer.

\subsection{Laser-induced electron diffraction}
\subsection{Femtosecond X-ray scattering}

\subsection{In the biological sciences}
Molecular movies are of course not only of interest in physics and chemistry as a means of probing fundamental processes, but also in the biological sciences where molecular structure play a crucial role in determining the function of biomolecules such as proteins. However, the molecules of interest there are much too large to be studied by any of the previous techniques. Thus molecular movies in the biological sciences tend to be annotated computer simulations amalgamated from multiple studies. That said, they are very impressive pieces of work.

A particularly impressive movie by \citet{Cheung12} showcases the process of RNA polymerase transcription and goes on for over six minutes.

\section{Previous attempts}

\subsection{Simplex algorithm}

\index{Simplex algorithm}
\index{Geometry reconstruction!Simplex algorithm}
\citet{Brichta09} proposed the reconstruction of small triatomic molecules using a simplex algorithm. Unfortunately they only report on the reconstruction of molecular structures based on simulated data for carbon dioxide and formaldehyde. I could not use this algorithm to find the geometries of CO2 or OCS from real data.

\subsection{Lookup table}
\index{Lookup table}
\index{Geometry reconstruction!Lookup table}
In this approach, many Coulomb explosions are simulated for a wide variety of structures, and the resulting momentum vectors from each simulation are stored. Thus you have a mapping from molecular structures to momentum vectors. To determine the structure belonging to a certain set of observed momentum vectors, you simply read the table in reverse. This approach is simple to implement, very quick by design, and front-loads the computation which may be desirable for large data sets. However, of course, it has an exponential time and space complexity $\mathcal{O}(e^{3N-6})$ where $N$ is the number of atoms.

\subsection{Convex optimization methods}
Much progress was made over the lookup table by treating the geometry reconstruction problem as a constrained nonlinear convex optimization such that MATLAB's \texttt{fmincon} function an be relied on. It relies on trust regions and uses an interior-point algorithm. This worked especially well in the case of triatomic molecules however four-atom systems proved incredibly difficult to tackle here. This was due to the exponential increase in the number of saddle points with dimensionality \marginpar{Recall that triatomic molecules have three degrees of freedom resulting in a problem of dimension 3 while four-atom systems have six. That is, $3N+6$ for an $N$-atom system.} making the problem highly non-convex and unsuitable for \texttt{fmincon}.

\section{Current attempt: MMMGRUBS}


\appendix
%\renewcommand{\thechapter}{\alph{chapter}}
\cleardoublepage

%\part{Appendix}
%\include{Chapters/Chapter0A}
\cleardoublepage%********************************************************************
% Bibliography
%*******************************************************
% work-around to have small caps also here in the headline
\manualmark
\markboth{\spacedlowsmallcaps{\bibname}}{\spacedlowsmallcaps{\bibname}} % work-around to have small caps also
%\phantomsection 
\refstepcounter{dummy}
\addtocontents{toc}{\protect\vspace{\beforebibskip}} % to have the bib a bit from the rest in the toc
\addcontentsline{toc}{chapter}{\tocEntry{\bibname}}
\label{app:bibliography}
\printbibliography

%\cleardoublepage%*******************************************************
% Declaration
%*******************************************************
\refstepcounter{dummy}
\pdfbookmark[0]{Declaration}{declaration}
\chapter*{Declaration}
\thispagestyle{empty}
Put your declaration here.
\bigskip
 
\noindent\textit{\myLocation, \myTime}

\smallskip

\begin{flushright}
    \begin{tabular}{m{5cm}}
        \\ \hline
        \centering\myName \\
    \end{tabular}
\end{flushright}

%\cleardoublepage\pagestyle{empty}

\hfill

\vfill


\pdfbookmark[0]{Colophon}{colophon}
\section*{Colophon}
This document was typeset using the typographical look-and-feel \texttt{classicthesis} developed by Andr\'e Miede. 
The style was inspired by Robert Bringhurst's seminal book on typography ``\emph{The Elements of Typographic Style}''. 
\texttt{classicthesis} is available for both \LaTeX\ and \mLyX: 
\begin{center}
\url{https://bitbucket.org/amiede/classicthesis/}
\end{center}
Happy users of \texttt{classicthesis} usually send a real postcard to the author, a collection of postcards received so far is featured here: 
\begin{center}
\url{http://postcards.miede.de/}
\end{center}
 
\bigskip

\noindent\finalVersionString

%Hermann Zapf's \emph{Palatino} and \emph{Euler} type faces (Type~1 PostScript fonts \emph{URW
%Palladio L} and \emph{FPL}) are used. The ``typewriter'' text is typeset in \emph{Bera Mono}, 
%originally developed by Bitstream, Inc. as ``Bitstream Vera''. (Type~1 PostScript fonts were made 
%available by Malte Rosenau and
%Ulrich Dirr.)

%\paragraph{note:} The custom size of the textblock was calculated
%using the directions given by Mr. Bringhurst (pages 26--29 and
%175/176). 10~pt Palatino needs  133.21~pt for the string
%``abcdefghijklmnopqrstuvwxyz''. This yields a good line length between
%24--26~pc (288--312~pt). Using a ``\emph{double square textblock}''
%with a 1:2 ratio this results in a textblock of 312:624~pt (which
%includes the headline in this design). A good alternative would be the
%``\emph{golden section textblock}'' with a ratio of 1:1.62, here
%312:505.44~pt. For comparison, \texttt{DIV9} of the \texttt{typearea}
%package results in a line length of 389~pt (32.4~pc), which is by far
%too long. However, this information will only be of interest for
%hardcore pseudo-typographers like me.%
%
%To make your own calculations, use the following commands and look up
%the corresponding lengths in the book:
%\begin{verbatim}
%    \settowidth{\abcd}{abcdefghijklmnopqrstuvwxyz}
%    \the\abcd\ % prints the value of the length
%\end{verbatim}
%Please see the file \texttt{classicthesis.sty} for some precalculated 
%values for Palatino and Minion.
%
%    \settowidth{\abcd}{abcdefghijklmnopqrstuvwxyz}
%    \the\abcd\ % prints the value of the length






\end{document}
